\documentclass[11pt]{article}
\usepackage[utf8]{inputenc}
\usepackage[spanish]{babel}
\usepackage{amsmath,amssymb,amsthm}
\usepackage{graphicx}
\usepackage[margin=2.5cm]{geometry}
\usepackage{enumitem}

\title{Semana 2 -- Ejercicios}
\author{Tutorías -- Introducción al Cálculo}
\date{10/08--14/08/2026}

\begin{document}
\maketitle

\section{Ejercicios de la semana}

\begin{enumerate}[label=\arabic*.]

\item Encuentre todas las raíces racionales de $P(x) = x^3 - 2x^2 - 25x + 50$
usando la Proposición de raíces racionales. (Respuesta: $2, 5, -5$.)

\item Encuentre las raíces racionales de $P(x) = 2x^3 + x^2 - 13x + 6$ y
factorice $P$ completamente usando el Teorema del Factor. (Respuesta:
raíces $\tfrac12, 2, -3$.)

\item Encuentre constantes $a,b \in \mathbb{R}$ tales que
$$x^3 + ax^2 + bx - 6 = (x-1)(x+2)(x-3) \quad \text{para todo } x.$$

\item Un polinomio $P$ de grado $3$ tiene exactamente las raíces $1, -1, 4$
y el coeficiente de $x^3$ es $2$. Determine $P(x)$ completamente,
justificando con el resultado que garantiza que un polinomio de grado $n$
queda determinado por $n$ raíces distintas y su coeficiente principal.

\item Sean $x,y,z \in \mathbb{R}$ con $x<y$ y $z<0$. Use los axiomas de orden
(y el hecho de que $-z>0$) para argumentar por qué se espera que
$xz > yz$ (es decir, que la desigualdad se invierte al multiplicar por un
número negativo).

\item Demuestre, usando los axiomas de orden, que si $x>0$ e $y>0$ entonces
$x+y>0$.

\item Demuestre que $\neg(x<y) \iff x\geq y$ usando el axioma de orden lineal
y la definición de $\geq$.

\item Sin usar la fórmula cuadrática, explique por qué $P(x)=x^2+4$ no puede
tener raíces racionales (ni reales), usando lo visto sobre positividad
($x^2\geq 0$).

\item Verifique si $\tfrac{3}{2}$ es raíz de $P(x) = 2x^3 - 3x^2 - 8x + 12$.
¿Es un candidato válido según la Proposición de raíces racionales?
(Respuesta: sí es raíz.)

\item Use los axiomas de orden para demostrar que si $0<x<y$ entonces
$x^2 < xy$ (ayuda: multiplique la desigualdad $x<y$ por $x$).

\item \textbf{(Nivel interrogación)} Un polinomio $P$ de grado $3$ tiene exactamente las raíces
racionales $1,-2,3$ y coeficiente principal $2$. Usando el Teorema 7.4 (y su corolario sobre
polinomios con la misma cantidad excesiva de raíces), determine $P(x)$ completamente y explique
rigurosamente por qué $P$ no puede tener una cuarta raíz distinta de las anteriores.

\item \textbf{(Nivel interrogación)} Demuestre rigurosamente, usando \textbf{solo} los axiomas de
orden, que si $0<a<b$ y $0<c<d$, entonces $ac<bd$. Indique explícitamente qué axioma utiliza en
cada paso.

\item \textbf{(Nivel interrogación)} Sea $P(x) = x^4 - 5x^3 + 5x^2 + 5x - 6$. Sabiendo que $1$ y
$-1$ son raíces de $P$, factorice $P(x)$ completamente (encuentre todas sus raíces) aplicando el
Teorema del Factor repetidamente.

\end{enumerate}

\end{document}
